% TOP_PROBLEM: {"id": 20001107, "title": "Dense nonsumsets and cyclic-interval square roots", "category_id": 2, "category": {"id": 2, "name": "combinatorics", "display_name": "Combinatorics", "description": "Counting problems, graph theory, discrete structures.", "slug": "combinatorics", "order_index": 2, "created_at": "2026-07-31T15:26:25.671Z"}, "status": "partially_solved", "problem_number": "AIM-COMBINATORICS-0232", "set_id": 14}
% FIRST_POSTED: 2026-10-01
% ENTRY_KIND: statement_only
% UnsolvedMath ID 20001107; source catalogue code AIM-COMBINATORICS-0232
% !TeX program = lualatex
% Definition and sources only; this file is not an LLM solution attempt.
\documentclass[11pt,a4paper]{article}
\usepackage[margin=25mm]{geometry}
\usepackage{fontspec}
\setmainfont{lmroman10-regular.otf}[BoldFont=lmroman10-bold.otf,
 ItalicFont=lmroman10-italic.otf,BoldItalicFont=lmroman10-bolditalic.otf]
\usepackage{amsmath,amssymb,mathtools}
\usepackage{unicode-math}
\setmathfont{latinmodern-math.otf}
\AtBeginDocument{\let\setminus\smallsetminus}
\usepackage{xurl,hyperref}
\hypersetup{colorlinks=true,urlcolor=blue}
\setcounter{secnumdepth}{0}
\setlength{\parindent}{0pt}
\setlength{\parskip}{5pt}
\setlength{\emergencystretch}{3em}
\begin{document}
\section*{Dense nonsumsets and cyclic-interval square roots}\label{20001107}
{\small UnsolvedMath ID 20001107 · AIM-COMBINATORICS-0232 · Combinatorics · AIM Workshop Problem Lists}
\subsection{Definitions and mathematical statement}
Let $p$ be an odd prime and regard $\mathbb F_p$ as an additive group.
For $B\subseteq\mathbb F_p$, define
$B+B=\{b_1+b_2:b_1,b_2\in B\}$, allowing $b_1=b_2$ and counting distinct sums only.
Call $S\subseteq\mathbb F_p$ a double sumset if there exists some $B\subseteq\mathbb F_p$ with $S=B+B$ exactly.
The set $B$ is not specified in advance and need not be a subset of $S$.
The empty set is a double sumset by taking $B=\varnothing$.

Define
\[
 M(p)=\max\{|S|:S\subseteq\mathbb F_p,
                   \text{no }B\subseteq\mathbb F_p\text{ has }B+B=S\},
 \qquad D(p)=p-M(p).
\]
Determine the asymptotic size of $M(p)$, or more informatively of its deficit $D(p)$, as $p$ ranges over primes tending to infinity.
Sharp bounds require both constructing large sets that fail to be double sumsets and proving that every set above a threshold is one.
The equality requirement is crucial: containment of $S$ in a double sumset would be trivial by using the whole field.
Likewise, representing $S$ as $B+C$ with two independently chosen sets is not the question.
The optimization is over all subsets, with no symmetry, interval, multiplicative, or random-model restrictions.
The source's historical bounds motivate the extremal function but are not taken as an assertion of the present best estimates.
\subsection{Short English statement}
How large can a subset of a prime field be while still failing to equal the sum of some set with itself?
\subsection{Sources}
\begin{itemize}
\item[S1] ULAM.AI and the UnsolvedMath Contributors, supplied problems.json, record 20001107 (AIM-COMBINATORICS-0232), AIM Workshop Problem Lists. Catalogue identity and source transcription; CC BY 4.0.
\url{https://huggingface.co/datasets/ulamai/UnsolvedMath}.
\item[S2] American Institute of Mathematics, Recent trends in additive combinatorics, item 2.7. Original workshop question underlying AIM-COMBINATORICS-0232.
\url{https://aimath.org/WWN/additivecomb/additivecomb.pdf}.
\item[S3] Noga Alon, Large sets in finite fields are sumsets. Extremal nonsumset context.
\url{https://web.math.princeton.edu/~nalon/PDFS/sumset.pdf}.
\end{itemize}
\end{document}
